%=============================================================
%  Module I.4 / L01 — Free Particle and Particle in a Box
%  Series I > Module I.4 > Lecture L01
%  Duration: 90 minutes  |  All tiers (T1–T5)
%  Input into master_document.tex OR compile standalone
%=============================================================
\renewcommand{\lectureCode}{I.4 / L01}

\section*{L01 — Free Particle and Particle in a Box}
\label{sec:L01}
\addcontentsline{toc}{section}{L01 — Free Particle and Particle in a Box}

%--- Lecture header ------------------------------------------
\begin{tcolorbox}[lecturebox, title={Module I.4 / L01 — Metadata}]
\begin{tabular}{@{}ll@{}}
\textbf{Series:}    & I — Introductory Quantum Mechanics \\
\textbf{Module:}    & I.4 — Paradigm Physical Systems \\
\textbf{Lecture:}   & L01 — Free Particle and Particle in a Box \\
\textbf{Duration:}  & 90 minutes \\
\textbf{Tiers:}     & T1 (HS) · T2 (BegUG) · T3 (AdvUG) · T4 (MSc) · T5 (PhD) \\
\textbf{Preceding:} & Module I.3 / L10 (Open quantum systems — last I.3 lecture) \\
\textbf{Following:} & I.4 / L02 — Simple Bloch Electron \\
\textbf{Prerequisites:} & Modules I.1–I.3 complete \\
\end{tabular}

\medskip
\textbf{Summary:}
This lecture introduces the two simplest exactly-solvable quantum systems: the free
particle (continuous spectrum, plane waves, wave packets) and the infinite square well
(discrete spectrum, standing waves). Together they establish the central dichotomy
between confined and unconfined quantum motion, and provide the first encounter with
density of states — a quantity that underpins statistical mechanics, condensed matter,
and quantum field theory. \textit{Why it matters:} every subsequent lecture in I.4
references one or both of these systems.
\end{tcolorbox}

%------------------------------------------------------------
\subsection*{L01.0 — Learning Objectives (per tier)}
\label{sec:L01.0}

\tiertag{tier1col}{Tier 1 — HS}\quad
Identify that a free particle carries definite momentum (\emph{plane wave});
describe why confining a particle in a box forces discrete energies;
interpret zero-point energy as a consequence of the uncertainty principle.

\tiertag{tier2col}{Tier 2 — BegUG}\quad
Apply the time-independent Schrödinger equation to free and box potentials;
compute $\psi_n$, $E_n$ for the infinite square well;
construct and interpret a Gaussian wave packet; calculate group velocity.

\tiertag{tier3col}{Tier 3 — AdvUG}\quad
Derive the density of states $g(E)$ in 1D, 2D, and 3D from $k$-space counting;
prove orthonormality and completeness of the square-well basis;
derive dispersion broadening of the Gaussian wave packet exactly.

\tiertag{tier4col}{Tier 4 — MSc}\quad
Analyse the finite square well and transfer matrix;
derive van Hove singularities from $\nabla_k E = 0$;
solve the Airy-function problem for a linear potential (triangular well).

\tiertag{tier5col}{Tier 5 — PhD}\quad
Discuss self-adjoint extensions of $\hat{p}$ on $[0,L]$ and their physical meaning;
formulate resonances and quasi-bound states via complex scaling;
connect the box density of states to the spectral zeta function.

%------------------------------------------------------------
\subsection*{L01.1 — Free Particle: Spectrum and Plane Waves}
\label{sec:L01.1}

\begin{tcolorbox}[lecturebox, title={L01.1 Main result}]
\textbf{Hamiltonian:}
\begin{equation}
  \hat{H} = \frac{\hat{p}^2}{2m}, \qquad
  \hat{H}\psi_k = E_k\psi_k, \qquad
  E_k = \frac{\hbar^2 k^2}{2m}.
  \label{eq:L01.1}
\end{equation}
\textbf{Plane-wave solutions} (box normalised, length $L$, periodic BCs):
\begin{equation}
  \psi_k(x) = \frac{1}{\sqrt{L}}\,e^{ikx}, \qquad
  k = \frac{2\pi n}{L},\; n\in\mathbb{Z}.
  \label{eq:L01.2}
\end{equation}
Continuum limit ($L\to\infty$): $\psi_k(x) = e^{ikx}/\sqrt{2\pi}$.
\end{tcolorbox}

\textbf{Physical interpretation.}
Plane waves carry definite momentum $p = \hbar k$ and are completely delocalised
($|\psi_k|^2 = \mathrm{const}$). They form an improper basis — normalisable only
in a distributional sense ($\langle k|k'\rangle = \delta(k-k')$).

\tiertag{tier5col}{Tier 5} \quad
The operator $\hat{p} = -i\hbar\partial_x$ on the half-line $[0,\infty)$ has
\emph{no} self-adjoint extension; on $[0,L]$ it has a one-parameter family
of self-adjoint extensions labelled by $e^{i\theta}$ (quasi-periodic BCs).

%------------------------------------------------------------
\subsection*{L01.2 — Gaussian Wave Packets}
\label{sec:L01.2}

\begin{tcolorbox}[lecturebox, title={L01.2 Main result}]
A wave packet is a superposition of plane waves:
\begin{equation}
  \psi(x,t) = \int \tilde{\phi}(k)\,e^{ikx - i\omega_k t}\,\frac{dk}{\sqrt{2\pi}},
  \qquad \omega_k = \frac{\hbar k^2}{2m}.
  \label{eq:L01.3}
\end{equation}
For a Gaussian envelope centred at $k_0$ with width $\sigma_k$:
\begin{equation}
  \tilde{\phi}(k) = \left(\frac{1}{2\pi\sigma_k^2}\right)^{1/4}
    e^{-(k-k_0)^2/4\sigma_k^2}.
  \label{eq:L01.4}
\end{equation}
\textbf{Group velocity:} $v_g = d\omega/dk\big|_{k_0} = \hbar k_0/m$.\\
\textbf{Phase velocity:} $v_p = \omega/k = \hbar k_0/2m = v_g/2$.\\
\textbf{Spreading:} $\Delta x(t) = \Delta x(0)\,\sqrt{1 + (t/t_*)^2}$ with
$t_* = 2m\Delta x(0)^2/\hbar$.
\end{tcolorbox}

\begin{example}[Electron wave packet on atomic scale]
Take $k_0 = 1\;\mathrm{nm}^{-1}$, $\Delta x(0) = 0.5\;\mathrm{nm}$.
Then $v_g = \hbar k_0/m_e \approx 1.15\times10^5\;\mathrm{m/s}$ and
$t_* \approx 4.8\;\mathrm{fs}$. An electron wave packet spreads to twice its
initial size in $\sim5$ femtoseconds — directly observable in attosecond experiments.
\end{example}

%------------------------------------------------------------
\subsection*{L01.3 — Infinite Square Well}
\label{sec:L01.3}

\begin{tcolorbox}[lecturebox, title={L01.3 Main result}]
Potential: $V(x) = 0$ for $0 < x < L$; $V = \infty$ otherwise.
\begin{align}
  \psi_n(x) &= \sqrt{\frac{2}{L}}\sin\!\left(\frac{n\pi x}{L}\right), \quad
  n = 1, 2, 3, \ldots \label{eq:L01.5}\\
  E_n &= \frac{n^2\pi^2\hbar^2}{2mL^2} = n^2 E_1,
  \qquad E_1 = \frac{\pi^2\hbar^2}{2mL^2}. \label{eq:L01.6}
\end{align}
\textbf{Orthonormality:} $\langle\psi_m|\psi_n\rangle = \delta_{mn}$.\\
\textbf{Completeness:} $\sum_n |\psi_n\rangle\langle\psi_n| = \hat{\mathbb{1}}$ on $[0,L]$.
\end{tcolorbox}

\textbf{Zero-point energy.}
Even the ground state has $E_1 > 0$. Physically: confining a particle to size $L$
gives $\Delta x \lesssim L$, so by the HUP $\Delta p \geq \hbar/2L$, and kinetic
energy $\langle\hat{p}^2\rangle/2m \geq \hbar^2/8mL^2$. No particle can be
\emph{at rest} inside a box.

\textbf{Expectation values:}
\begin{equation}
  \langle x\rangle_n = \frac{L}{2}, \quad
  \langle x^2\rangle_n = L^2\!\left(\frac{1}{3} - \frac{1}{2n^2\pi^2}\right), \quad
  \langle p\rangle_n = 0, \quad
  \langle p^2\rangle_n = \frac{n^2\pi^2\hbar^2}{L^2}.
  \label{eq:L01.7}
\end{equation}

%------------------------------------------------------------
\subsection*{L01.4 — Density of States in 1D, 2D, 3D}
\label{sec:L01.4}

\begin{tcolorbox}[lecturebox, title={L01.4 Main result}]
Density of states $g(E)$ = number of states per unit energy per unit volume:
\begin{align}
  g_{1D}(E) &= \frac{1}{\pi\hbar}\sqrt{\frac{m}{2E}} \label{eq:L01.8}\\
  g_{2D}(E) &= \frac{m}{\pi\hbar^2} \quad (\text{constant!}) \label{eq:L01.9}\\
  g_{3D}(E) &= \frac{(2m)^{3/2}}{2\pi^2\hbar^3}\sqrt{E} \label{eq:L01.10}
\end{align}
\end{tcolorbox}

\textbf{Derivation sketch (3D):}
Count $k$-states in shell $[k, k+dk]$: $dn = V\cdot 4\pi k^2 dk/(2\pi)^3$.
Use $E = \hbar^2k^2/2m$ to eliminate $k$, divide by $V$ to get $g(E)$.

\tiertag{tier4col}{Tier 4}\quad
\textbf{Van Hove singularities:} whenever $\nabla_k E = 0$, $g(E)$ diverges
(1D: $E^{-1/2}$ divergence; 2D: logarithmic jump; 3D: saddle-point discontinuity in slope).
These are directly observable in angle-resolved photoemission spectra.

%------------------------------------------------------------
\subsection*{L01.5 — Connections}
\label{sec:L01.5}

\textbf{Backward:} I.3/L01 (time evolution operator $\hat{U}(t) = e^{-i\hat{H}t/\hbar}$) —
the wave packet propagation in L01.2 is a direct application.

\textbf{Forward:} I.4/L02 — the Bloch electron in a periodic potential generalises
the free particle; the free particle DOS becomes the band DOS.
I.4/L03 — the harmonic oscillator also has equally spaced levels, but in energy, not $k$.

\textbf{Cross-module:} II.4/L01 (Fock space) — the box modes become modes of the
quantised EM field; creation/annihilation operators replace $\hat{a}$, $\hat{a}^\dagger$.

%=============================================================
%  ASSESSMENT BUNDLE — L01
%=============================================================
\newpage
\section*{L01 — Assessment Artifact Bundle}

%------------------------------------------------------------
\subsection*{[A] SCQU — Simple Concept Questions (UG, ×10)}
\begin{tcolorbox}[ugconceptbox, title={SCQU · L01 · Tiers 1–3 · 15–20 min}]
\begin{enumerate}[label=\textbf{Q\arabic*.}]
  \item \textit{(MC)} The group velocity of a free-particle wave packet is:
    (a) $\hbar k/m$ \quad (b) $\hbar k/2m$ \quad (c) $\hbar k^2/2m$ \quad (d) $k/\hbar$.
  \item \textit{(MC)} In an infinite square well of width $L$, the ground-state energy scales as:
    (a) $L^{-1}$ \quad (b) $L^{-2}$ \quad (c) $L^{1/2}$ \quad (d) $L^{-3}$.
  \item \textit{(T/F + justify)} The phase velocity and group velocity of a free particle are equal.
  \item \textit{(T/F + justify)} A particle in the $n=2$ state of an infinite well has a node at $x = L/2$.
  \item \textit{(T/F + justify)} The density of states in 2D is constant (independent of $E$).
  \item \textit{(T/F + justify)} Zero-point energy violates conservation of energy.
  \item \textit{(Short)} Explain physically why $\langle p \rangle = 0$ for every stationary state of the infinite well.
  \item \textit{(Short)} What happens to the energy eigenvalues of the box if $L$ is doubled?
  \item \textit{(MC)} A Gaussian wave packet with $\Delta x(0) = \sigma$ has uncertainty-product $\Delta x \cdot \Delta p$ at $t=0$ equal to:
    (a) $\hbar/2$ \quad (b) $\hbar$ \quad (c) $\hbar\sigma$ \quad (d) depends on $k_0$.
  \item \textit{(T/F + justify)} In 3D, $g(E) \propto E^{1/2}$ so states accumulate at higher energies.
\end{enumerate}
\textit{Rubric: 2 pts each (1.5 correct answer + 0.5 justification). Total: 20 pts.}
\end{tcolorbox}

%------------------------------------------------------------
\subsection*{[B] CCQU — Complex Concept Questions (UG, ×10)}
\begin{tcolorbox}[ugconceptbox, title={CCQU · L01 · Tiers 2–3 · 25–35 min}]
\begin{enumerate}[label=\textbf{Q\arabic*.}]
  \item Explain why a plane wave $e^{ikx}$ cannot be normalised on $\mathbb{R}$, but can on $[0,L]$ with periodic BCs.
  \item Why does confining a particle to a smaller box \emph{increase} its zero-point energy? Connect to the uncertainty principle.
  \item What would change if we allowed $n=0$ as a quantum number in the infinite square well? Is $\psi_0 = 0$ physically acceptable?
  \item Identify the flaw: ``Since $\langle p \rangle = 0$ for all box eigenstates, the particle is at rest.''
  \item A classically confined particle bounces back and forth with constant speed. Explain how $|\psi_n|^2 \to \mathrm{const}$ for large $n$ (correspondence principle) and why this makes sense.
  \item What would change in the density of states formula if the particle were a photon ($E = \hbar ck$, linear dispersion)?
  \item Explain why $g_{2D}(E) = \mathrm{const}$ while $g_{1D}(E)$ diverges at $E=0$ and $g_{3D}(E) \to 0$ as $E\to0$.
  \item ``A wave packet is just a superposition, so it should spread indefinitely.'' What is correct and what is missing in this statement?
  \item How does the orthonormality $\langle\psi_m|\psi_n\rangle = \delta_{mn}$ follow from the boundary conditions, not just from the Hamiltonian being Hermitian?
  \item The 3D density of states $g(E) \propto E^{1/2}$ plays a central role in Fermi–Dirac statistics. Explain intuitively why higher energies have more available states.
\end{enumerate}
\textit{Rubric: 2 pts each (1 key point + 1 argument). Total: 20 pts.}
\end{tcolorbox}

%------------------------------------------------------------
\subsection*{[C] SPSU — Simple Problem Set (UG, ×10)}
\begin{tcolorbox}[ugprobbox, title={SPSU · L01 · Tiers 1–3 · 60–90 min}]

\textbf{Part A: Mechanics (4 problems)}
\begin{enumerate}[label=\textbf{P\arabic*.}]
  \item For an infinite square well of width $L=1\;\mathrm{nm}$, compute $E_1$, $E_2$, $E_3$
    in eV for an electron. What photon wavelength corresponds to the $n=1\to2$ transition?
  \item Normalise the Gaussian wave packet $\tilde{\phi}(k) = A e^{-(k-k_0)^2/4\sigma_k^2}$ and
    find the position-space form $\psi(x,0)$ by Fourier transform.
  \item Compute $\langle x\rangle$, $\langle x^2\rangle$, $\langle p^2\rangle$ for the $n=3$
    state of the infinite square well and verify $E_3 = \langle p^2\rangle/2m$.
  \item Show that the infinite-well eigenfunctions satisfy
    $\int_0^L \psi_m^*(x)\psi_n(x)\,dx = \delta_{mn}$ explicitly for $(m,n) = (1,2)$.
\end{enumerate}

\textbf{Part B: Interpretation (3 problems)}
\begin{enumerate}[label=\textbf{P\arabic*.},resume]
  \item The ground-state energy of a proton in a nucleus ($L \sim 1\;\mathrm{fm}$) is roughly
    what fraction of its rest-mass energy $m_p c^2 = 938\;\mathrm{MeV}$?
    Comment on whether non-relativistic QM is adequate.
  \item A state is prepared as $\ket{\psi} = \frac{1}{\sqrt{2}}(\ket{1}+\ket{2})$ in the box.
    Compute $\langle E\rangle$, $\langle x\rangle(t)$, and $|\langle x\rangle(t)|$.
  \item Sketch $|\psi_n(x)|^2$ for $n=1,4,10$. Describe the approach to the classical
    uniform distribution.
\end{enumerate}

\textbf{Part C: Translation (3 problems)}
\begin{enumerate}[label=\textbf{P\arabic*.},resume]
  \item The density of states in 3D is $g(E) = V(2m)^{3/2}E^{1/2}/2\pi^2\hbar^3$.
    Derive the total number of states with energy $\leq E_F$ (Fermi energy) for $N$ free electrons.
  \item Express the group velocity $v_g = \hbar k_0/m$ in terms of the de Broglie wavelength $\lambda$ and period $T$.
  \item In 2D, the DOS is constant $g_{2D} = m/\pi\hbar^2$. Use this to find the number of
    electrons per unit area at $T=0$ for a 2D electron gas with Fermi energy $E_F$.
\end{enumerate}
\textit{Rubric: Result 60\%, Method 30\%, Interpretation 10\%.}
\end{tcolorbox}

%------------------------------------------------------------
\subsection*{[D] CPSU — Complex Problem Set (UG, ×5, 4-layer format)}
\begin{tcolorbox}[ugprobbox, title={CPSU · L01 · Tier 3 · 3–6 hrs total}]

\textbf{Layer 1 — Problem Statements}
\begin{enumerate}[label=\textbf{Ex\arabic*.}]
  \item \textbf{Wave-packet spreading.}
    (a) Show that for any initial state $\psi(x,0)$ in a free-particle Hilbert space,
    $\Delta x(t)^2 = \Delta x(0)^2 + (\Delta p/m)^2 t^2 + C\cdot t$ for some constant $C$.
    (b) For a Gaussian packet, show $C=0$ and derive $\Delta x(t)$ exactly.
    (c) Interpret the result: what is $t_*$ physically?

  \item \textbf{Completeness and Parseval.}
    (a) Using completeness of the infinite-well basis, prove Parseval's identity:
    $\int_0^L |\psi(x)|^2\,dx = \sum_n |c_n|^2$ where $c_n = \langle\psi_n|\psi\rangle$.
    (b) Expand $f(x) = x(L-x)$ in the infinite-well basis and verify Parseval numerically
    (truncate at $n=20$).

  \item \textbf{3D box degeneracy.}
    (a) For a 3D cubic box of side $L$, write the energy levels $E_{n_x,n_y,n_z}$.
    (b) List all distinct $(n_x,n_y,n_z)$ triples with energy $E \leq 6E_1$
    and determine the degeneracy of each level.
    (c) Show that the level $E = 9E_1$ has degeneracy 3 but $E = 14E_1$ has degeneracy 6.

  \item \textbf{Uncertainty principle from box eigenstates.}
    (a) Compute $\Delta x \cdot \Delta p$ for the $n$-th infinite-well eigenstate.
    (b) Show it exceeds $\hbar/2$ for all $n\geq1$.
    (c) Find the limit as $n\to\infty$ and interpret using the correspondence principle.

  \item \textbf{Finite versus infinite well preview.}
    (a) In a finite square well of depth $V_0$ and width $a$, write the boundary conditions
    linking the interior and exterior wave functions.
    (b) Show that there is always at least one bound state (no matter how small $V_0 a^2$
    is in 1D) using the graphical transcendental equation.
    (c) Estimate the number of bound states for $V_0 = 10\;\mathrm{eV}$, $a = 1\;\mathrm{nm}$, $m = m_e$.
\end{enumerate}
\textit{Rubric: Assumptions 15\%, Proof structure 45\%, Algebra 25\%, Insight 15\%.}\\
\textit{See} \texttt{ex01\_hints.tex}\textit{, }\texttt{ex01\_adv\_hints.tex}\textit{, }
\texttt{ex01\_solution.tex} \textit{for Layers 2–4.}
\end{tcolorbox}

%------------------------------------------------------------
\subsection*{[E] SPJU — Simple Project (UG)}
\begin{tcolorbox}[ugprobbox, title={SPJU · L01 · Tiers 2–3 · 2–4 hrs}]
\textbf{Project: Gaussian Wave Packet Visualiser}

\textit{Goal:} Write a Python script (or Jupyter notebook) that animates $|\psi(x,t)|^2$
for a Gaussian wave packet in free space.

\textit{Inputs:}
\begin{itemize}[topsep=2pt]
  \item Initial position $x_0$, width $\sigma_x$, central momentum $k_0$
  \item Time range $[0, T_{\max}]$ and number of frames
\end{itemize}

\textit{Required output:}
\begin{enumerate}[topsep=2pt]
  \item Animated plot of $|\psi(x,t)|^2$ showing spreading and translation
  \item Plot of $\Delta x(t)$ vs $t$ with theoretical curve $\Delta x(0)\sqrt{1+(t/t_*)^2}$
  \item Numerical verification that $\int|\psi|^2\,dx = 1$ at each timestep (to $<0.1\%$ error)
\end{enumerate}

\textit{Deliverables:} Code + README + 1–2 page report explaining:
(1) how you implemented the FFT-based propagation, (2) the physical interpretation of $t_*$,
(3) one unexpected observation from your animation.

\textit{Rubric:} Correctness 40\%, Verification tests 25\%, Clarity 25\%, Reproducibility 10\%.
\end{tcolorbox}

%------------------------------------------------------------
\subsection*{[F] CPJU — Complex Project (UG)}
\begin{tcolorbox}[ugprobbox, title={CPJU · L01 · Tier 3 · 8–15 hrs}]
\textbf{Project: Van Hove Singularities in 1D, 2D, and 3D}

\textit{Goal:} Investigate the analytical structure of $g(E)$ in 1, 2, and 3 dimensions
for a free particle, and connect to experimental observables.

\textit{Deliverables:} 5–8 page report + code appendix addressing:
\begin{enumerate}[topsep=2pt]
  \item Derive $g_d(E)$ for $d = 1, 2, 3$ from first principles (no looking up the formula).
  \item Show that the 1D DOS diverges as $E^{-1/2}$ near $E=0$. Explain what this means
    for the low-temperature heat capacity of a 1D electron gas.
  \item In 2D, the DOS is constant: numerically verify this by computing $g(E)$ via
    counting states in a large 2D box and plotting vs $E$.
  \item For a 1D tight-binding chain $E(k) = -2t\cos(ka)$, compute $g(E)$ analytically
    and show it diverges at band edges (van Hove singularities).
  \item Compare your 1D tight-binding $g(E)$ to a real material measurement (e.g., DOS
    of a carbon nanotube from literature) and comment on qualitative agreement.
\end{enumerate}
\textit{Rubric:} Technical 30\%, Synthesis 35\%, Depth 25\%, Originality 10\%.
\end{tcolorbox}

%------------------------------------------------------------
\subsection*{[G] SCQG — Simple Concept Questions (Grad, ×10)}
\begin{tcolorbox}[gradconceptbox, title={SCQG · L01 · Tiers 4–5 · 20–25 min}]
\begin{enumerate}[label=\textbf{Q\arabic*.}]
  \item State the domain of $\hat{p} = -i\hbar\partial_x$ as a self-adjoint operator on $L^2[0,L]$.
  \item \textit{(T/F + proof sketch)} The operator $\hat{p}^2$ is self-adjoint on $L^2(\mathbb{R})$ with its natural domain.
  \item Identify the precise error: ``The plane-wave states $\{e^{ikx}\}$ form an orthonormal basis of $L^2(\mathbb{R})$.''
  \item \textit{(Short)} In what sense does $\delta(x-x_0)$ lie in the rigged Hilbert space $\Phi^* \supset L^2 \supset \Phi$?
  \item \textit{(Short)} State the spectral theorem for $\hat{H} = -\hbar^2\partial_x^2/2m$ on $[0,L]$ with Dirichlet BCs.
  \item \textit{(T/F + proof sketch)} The resolvent $(\hat{H}-z)^{-1}$ is a compact operator on $L^2[0,L]$ for $z\notin\sigma(\hat{H})$.
  \item What is the precise statement of Parseval's theorem in $L^2[0,L]$ and how does it follow from the spectral theorem?
  \item \textit{(Identify error)} ``The ground-state energy $E_1 \to 0$ as $L\to\infty$, so the particle becomes free.''
  \item \textit{(Short)} Define the distributional derivative and explain why $\hat{p}$ acts on wave functions in $H^1[0,L]$.
  \item \textit{(T/F + proof sketch)} Van Hove singularities in the DOS are topologically protected in the sense that they cannot be removed by small perturbations of the dispersion relation.
\end{enumerate}
\textit{Rubric: 2 pts each (1.5 correct answer + 0.5 justification). Total: 20 pts.}
\end{tcolorbox}

%------------------------------------------------------------
\subsection*{[H] CCQG — Complex Concept Questions (Grad, ×10)}
\begin{tcolorbox}[gradconceptbox, title={CCQG · L01 · Tiers 4–5 · 35–50 min}]
\begin{enumerate}[label=\textbf{Q\arabic*.}]
  \item Distinguish the notions of ``Hermitian,'' ``symmetric,'' and ``self-adjoint'' for $\hat{p}$ on $[0,L]$. Which boundary conditions make $\hat{p}$ self-adjoint?
  \item What fails in the spectral expansion $\psi = \sum_n c_n\psi_n$ if the Hamiltonian is not self-adjoint?
  \item Connect the van Hove singularities of the DOS to the topology of the Fermi surface via Morse theory.
  \item Explain why the density of states in 3D is $g(E)\propto E^{1/2}$ and how this changes for a relativistic dispersion $E = \hbar ck$.
  \item ``Particle in a box is trivial.'' Argue against this by identifying three non-trivial mathematical structures it encodes.
  \item What is the physical content of the spectral zeta function $\zeta_H(s) = \sum_n E_n^{-s}$ for the box Hamiltonian?
  \item Distinguish the spectral theory of the box Hamiltonian on $[0,L]$ (compact, discrete spectrum) from the free Hamiltonian on $\mathbb{R}$ (non-compact, continuous spectrum).
  \item Connect the Gaussian wave packet spreading formula to the Heisenberg time–energy uncertainty relation.
  \item What fails in the derivation of group velocity $v_g = d\omega/dk$ for a highly non-linear dispersion relation?
  \item How does the DOS change if we go from spinless particles to spin-$1/2$ fermions?
\end{enumerate}
\textit{Rubric: 2 pts each (1 key point + 1 argument). Total: 20 pts.}
\end{tcolorbox}

%------------------------------------------------------------
\subsection*{[I] SPSG — Simple Problem Set (Grad, ×10)}
\begin{tcolorbox}[gradprobbox, title={SPSG · L01 · Tiers 4–5 · 60–90 min}]
\textbf{Part A: Operator mechanics (4)}
\begin{enumerate}[label=\textbf{P\arabic*.}]
  \item Show that $\hat{p}$ on $H^1_0[0,L]$ (Dirichlet BCs) is \emph{symmetric} but not self-adjoint. Find its deficiency indices.
  \item Using the spectral theorem, express the propagator $e^{-i\hat{H}t/\hbar}$ as a sum over energy eigenstates and verify it maps $L^2[0,L]$ to $L^2[0,L]$ unitarily.
  \item The resolvent of the box Hamiltonian $G(x,y;z) = \langle x|(\hat{H}-z)^{-1}|y\rangle$ is a Green's function. Compute it explicitly for $z\notin\sigma(\hat{H})$.
  \item For the free Hamiltonian on $\mathbb{R}$, compute the spectral measure $dE_\lambda$ (i.e., the projector onto states with energy $\leq\lambda$) explicitly in position space.
\end{enumerate}
\textbf{Part B: Rigorous translation (3)}
\begin{enumerate}[label=\textbf{P\arabic*.},resume]
  \item Prove that the heat kernel $K(x,y;t) = \sum_n \psi_n(x)\psi_n^*(y)e^{-E_nt/\hbar}$ satisfies the heat equation $-\hbar\partial_t K = \hat{H}K$ with initial condition $K(x,y;0) = \delta(x-y)$.
  \item Show that $g(E)$ for the 3D box can be derived from the heat kernel trace $Z(\beta) = \mathrm{Tr}(e^{-\beta\hat{H}})$ via an inverse Laplace transform.
  \item Prove Weyl's law: the eigenvalue counting function $N(E) = \#\{n: E_n\leq E\}$ for the $d$-dimensional box satisfies $N(E) \sim V(2mE)^{d/2}/(\Gamma(d/2+1)(2\pi\hbar)^d)$.
\end{enumerate}
\textbf{Part C: Symmetry \& invariance (3)}
\begin{enumerate}[label=\textbf{P\arabic*.},resume]
  \item Under the dilation $x\to\lambda x$, $L\to\lambda L$, how does $E_n$ scale? Relate this to the virial theorem $2\langle T\rangle = n\langle V\rangle$ (here $n=2$ for the kinetic energy).
  \item Parity: the operator $\hat{\Pi}$ maps $x\to L-x$ in the shifted box $[-L/2,L/2]$. Identify the parity of each $\psi_n$ and verify $[\hat{H},\hat{\Pi}]=0$.
  \item Time-reversal: $\hat{T}\psi(x) = \psi^*(x)$. Show that $\hat{T}$ is antiunitary, maps $\psi_n\to\psi_n^*=\psi_n$ (real), and commutes with $\hat{H}$.
\end{enumerate}
\textit{Rubric: Result + conditions 50\%, Method 35\%, Meaning 15\%.}
\end{tcolorbox}

%------------------------------------------------------------
\subsection*{[J] CPSG — Complex Problem Set (Grad, ×5)}
\begin{tcolorbox}[gradprobbox, title={CPSG · L01 · Tiers 4–5 · 5–10 hrs}]
\begin{enumerate}[label=\textbf{Ex\arabic*.}]
  \item \textbf{Self-adjoint extensions of $\hat{p}$ on $[0,L]$.}
    Classify all self-adjoint extensions of $\hat{p}$ using von Neumann's deficiency index theory. Show that each extension corresponds to a quasi-periodic boundary condition $\psi(L) = e^{i\theta}\psi(0)$, $\theta\in[0,2\pi)$. Find the spectrum for each $\theta$.
  \item \textbf{Aharonov–Bohm ring.}
    A particle on a ring of circumference $L$ with a magnetic flux $\Phi$ threading the ring has Hamiltonian $\hat{H} = (\hat{p}-eA)^2/2m$ with $A = \Phi/L$. (a) Find the energy spectrum. (b) Show the spectrum is periodic in $\Phi$ with period $hc/e$ (flux quantum). (c) Connect to the self-adjoint extension parameter $\theta = 2\pi\Phi/\Phi_0$.
  \item \textbf{Weyl's law for the 3D box.}
    Prove that the eigenvalue counting function $N(E)$ satisfies $N(E) = V(2mE)^{3/2}/(6\pi^2\hbar^3) + O(E)$ as $E\to\infty$. The $O(E)$ correction comes from surface states — estimate it.
  \item \textbf{Spectral zeta function.}
    For the 1D box, compute $\zeta_H(s) = \sum_{n=1}^\infty E_n^{-s}$ in closed form. Show it equals $(mL^2/\pi^2\hbar^2)^s \zeta_R(2s)$ where $\zeta_R$ is the Riemann zeta function. Find the poles and their physical meaning.
  \item \textbf{Wave packet on a lattice.}
    Replace the free dispersion $E(k)=\hbar^2k^2/2m$ by the tight-binding dispersion $E(k) = -2t\cos(ka)$. (a) Compute the group velocity and effective mass. (b) Show a Gaussian wave packet centred at $k_0 = \pi/2a$ eventually breaks up (Zener tunnelling analogy). (c) Quantify the spreading rate vs. the free-particle case.
\end{enumerate}
\textit{See} \texttt{ex\_grad01\_hints.tex} \textit{etc.\ for Layers 2–4.}
\end{tcolorbox}

%------------------------------------------------------------
\subsection*{[K] SPJG — Simple Projects (Grad, ×5)}
\begin{tcolorbox}[gradprobbox, title={SPJG · L01 · Tiers 4–5 · 3–8 hrs each}]
\begin{enumerate}[label=\textbf{PG\arabic*.}]
  \item \textbf{Green's function of the box.} Compute $G(x,y;E) = \langle x|(\hat{H}-E)^{-1}|y\rangle$ analytically and numerically. Plot it and identify the poles. Deliverable: code + 2-page tech note.
  \item \textbf{Spectral measure of the free particle.} Numerically compute the spectral projector $E(\lambda) = \int_{-\infty}^\lambda |k\rangle\langle k|\,dk$ and verify it is a projection-valued measure. Deliverable: 2-page tech note + code.
  \item \textbf{DOS from the heat kernel.} Numerically extract $g(E)$ from $Z(\beta) = \mathrm{Tr}(e^{-\beta\hat{H}})$ via numerical inverse Laplace transform and compare to the analytical formula. Deliverable: code + 2-page tech note.
  \item \textbf{Quasi-periodic spectrum.} Implement the family of self-adjoint extensions of $\hat{p}$ on $[0,L]$, compute the spectrum as a function of $\theta$, and plot the Hofstadter-like butterfly. Deliverable: code + 2-page note.
  \item \textbf{Wave-packet in tight-binding chain.} Simulate the time evolution of a Gaussian wave packet in a 1D tight-binding chain. Compare spreading to free-particle case. Investigate Bloch oscillations under a constant force field. Deliverable: animated figure + 3-page note.
\end{enumerate}
\end{tcolorbox}

%------------------------------------------------------------
\subsection*{[L] CPJG — Complex Projects (Grad, ×5)}
\begin{tcolorbox}[gradprobbox, title={CPJG · L01 · Tiers 4–5 · 10–20 hrs each}]
\begin{enumerate}[label=\textbf{PG\arabic*.}]
  \item \textbf{Casimir energy of the box.} Using the spectral zeta function, regularise the zero-point energy $E_0 = \frac{1}{2}\sum_n E_n$ of the 1D box and compute the Casimir force $F = -\partial E_0/\partial L$. Deliverable: 6–10 page report.
  \item \textbf{Quantum chaos in the 2D box.} Numerically study the level spacing statistics of the 2D rectangular box as a function of the aspect ratio $L_x/L_y$. Compare to Poisson (integrable) statistics. Deliverable: 6–10 page report.
  \item \textbf{Transfer matrix and photonic crystals.} Extend the transfer matrix method from 1D potential barriers to a periodic sequence of layers (photonic crystal). Compute transmission bands and gaps. Compare to Bloch's theorem. Deliverable: 6–10 page report.
  \item \textbf{Wigner function of the infinite well.} Compute the Wigner function $W(x,p;t)$ for a superposition state in the infinite well. Visualise the non-classical interference fringes. Investigate the semiclassical limit $n\gg1$. Deliverable: 6–10 page report with figures.
  \item \textbf{Resonances from complex scaling.} Use the method of complex scaling ($x\to xe^{i\theta}$) to compute the energy and decay rate of a quasi-bound state in a potential well with a barrier. Compare to WKB tunnelling rate. Deliverable: 6–10 page report.
\end{enumerate}
\end{tcolorbox}

%------------------------------------------------------------
\subsection*{[M] WRQ — Well-Defined Research Questions (×5)}
\begin{tcolorbox}[researchbox, title={WRQ · L01 · Tiers 4–5 · 5–10 hrs each}]
\begin{enumerate}[label=\textbf{R\arabic*.}]
  \item \textbf{Particle-in-a-box as a model for conjugated molecules.}
    The free-electron model (PIB) predicts UV absorption wavelengths of linear polyenes.
    How accurate is it for butadiene, hexatriene, and $\beta$-carotene? What corrections
    improve the agreement? \textit{Deliverable: 2–4 pg memo with calculations and 5–10 references.}
  \item \textbf{van Hove singularities in carbon nanotubes.}
    Measured STM conductance spectra of single-walled carbon nanotubes show sharp peaks.
    Identify these as 1D van Hove singularities. Quantify the agreement between the free-electron
    model DOS and experimental data for a (10,0) nanotube.
  \item \textbf{Casimir effect between conducting plates.}
    Derive the Casimir force $F = -\pi^2\hbar c/(240 L^4) \cdot A$ between two parallel
    perfectly conducting plates separated by $L$, starting from the 3D electromagnetic
    box DOS. What is the precision of current experimental tests?
  \item \textbf{Wave-packet spreading in Bose–Einstein condensates.}
    After a BEC is released from a harmonic trap, it expands as a wave packet.
    How well does the free-particle spreading formula apply? What corrections arise from
    interactions (mean-field repulsion)? Reference at least two experiments.
  \item \textbf{Quantum confinement in semiconductor quantum dots.}
    The PIB model predicts that confining electrons to a nanoscale box blue-shifts the
    optical absorption. Quantitatively compare the PIB prediction to experimental absorption
    spectra of CdSe quantum dots of varying radius.
\end{enumerate}
\end{tcolorbox}

%------------------------------------------------------------
\subsection*{[N] ORQ — Open-Ended Research Questions (×5)}
\begin{tcolorbox}[researchbox, title={ORQ · L01 · Tiers 4–5 · 10–20 hrs each}]
\begin{enumerate}[label=\textbf{O\arabic*.}]
  \item \textbf{What is the ``correct'' quantisation of a particle on a manifold?}
    The particle-in-a-box problem on a curved manifold (e.g., a sphere, torus, or hyperbolic
    surface) raises the ``operator-ordering ambiguity.'' What are the competing proposals
    (Laplace–Beltrami, conformal, ...) and is there a physical way to distinguish them?
  \item \textbf{Is there a universal lower bound on wave-packet spreading?}
    The Gaussian wave packet saturates $\Delta x(0)\Delta p = \hbar/2$. Are there states
    that spread slower than Gaussian in free space? What does the quantum speed limit say?
  \item \textbf{Can the infinite square well be realised exactly?}
    Infinite potentials are unphysical. What do real semiconductor quantum wells achieve?
    How are energy levels modified by finite barrier height, effective mass discontinuity,
    and band non-parabolicity?
  \item \textbf{How does the density of states change in fractal and disordered geometries?}
    For a particle in a Sierpinski triangle or a disordered potential, $g(E)$ no longer
    follows power laws. What replaces the Weyl law? What are spectral dimension and fracton
    excitations?
  \item \textbf{What is the connection between the spectral zeta function of the box and the Riemann hypothesis?}
    Propose a physical system whose energy spectrum encodes the zeros of the Riemann zeta
    function (Berry–Keating conjecture). Summarise current mathematical and physical
    progress toward a proof.
\end{enumerate}
\end{tcolorbox}

%------------------------------------------------------------
\subsection*{Reading and Resources — L01}
\begin{tcolorbox}[goldbox]
\textbf{Primary (all tiers):}
Griffiths, \textit{Introduction to Quantum Mechanics}, Ch.~2 (\S2.1--2.3).
Sakurai \& Napolitano, \textit{Modern QM}, Ch.~2.4 (wave packets), Appendix~A.

\textbf{Core (Tiers 2–3):}
Cohen-Tannoudji et al., \textit{QM} Vol.~I, Complement~C$_\mathrm{I}$ (density of states).

\textbf{Graduate (Tiers 4–5):}
Reed \& Simon, \textit{Methods of Mathematical Physics}, Vol.~II (self-adjoint extensions).
Ashcroft \& Mermin, \textit{Solid State Physics}, Ch.~2 (Fermi sphere, DOS).

\textbf{Research:}
M.V.\ Berry \& M.\ Tabor, ``Level clustering in the regular spectrum,'' \textit{Proc.\ R.\ Soc.\ A} 356 (1977) — for ORQ Q2 on spectral statistics.
\end{tcolorbox}
